\documentclass[11pt]{article}
\usepackage[letterpaper,margin=0.92in]{geometry}
\usepackage[T1]{fontenc}
\usepackage{lmodern,microtype,amsmath,amssymb,booktabs,xcolor,fancyhdr}
\usepackage[colorlinks=true,linkcolor=ink,urlcolor=accent]{hyperref}
\definecolor{ink}{HTML}{182735}
\definecolor{accent}{HTML}{16656B}
\hypersetup{pdftitle={Five worked examples from a census of Collatz-like systems}}
\pagestyle{fancy}
\fancyhf{}
\fancyfoot[L]{\small\sffamily\color{accent} Collatz-like systems census}
\fancyfoot[R]{\small\sffamily\color{ink}\thepage}
\renewcommand{\headrulewidth}{0pt}
\setlength{\parindent}{0pt}
\setlength{\parskip}{5pt}
\newcommand{\tier}[2]{%
  \section*{\color{accent}\sffamily Tier #1 \quad #2}%
}
\newcommand{\record}[1]{{\small\sffamily\color{accent} Census record \quad\texttt{#1}}\par}
\begin{document}
\color{ink}
{\LARGE\sffamily\bfseries Five worked examples from a census\par
of Collatz-like systems\par}
\vspace{5pt}
{\small\sffamily\color{accent} TWO-ROUND CENSUS \quad / \quad RESEARCH NOTE}
\vspace{5pt}

\textbf{Abstract.} A two-round census ranked 179 questions about 112 Collatz-like
systems. The five examples below illustrate the methods that resolved questions
at tiers 0--4. The tier counts, in order from 0 to 5, are 43, 54, 12, 2, 2 and 66.
The census ranks the question and the form of its answer. The same map can have
a directly computable seed orbit and a global question that requires proof.
The census was written and cross-checked by agents and has had no human review.

\tier{0}{A single orbit}
\record{par\_1\_1\_s27}
Let \(T(n)=n/2\) for even \(n\) and \(T(n)=n+1\) for odd \(n\).
Which cycle does the orbit of \(27\) reach? Direct evaluation gives
\[
27\to28\to14\to7\to8\to4\to2\to1\to2.
\]
The repeated value \(2\) certifies that the orbit enters the cycle \((1,2)\).

\tier{1}{Descent for every starting value}
\record{par\_1\_1\_s27}
For the same map, does every positive orbit reach \((1,2)\)?
Every even \(n>2\) decreases in one step, and every odd \(n\ge3\)
decreases in two steps, since \(T^2(n)=(n+1)/2<n\).
Repeating these decreasing blocks must reach \(1\) or \(2\), since the
block endpoints are positive integers.

\tier{2}{Congruence classes and a finite table}
\record{up\_qr7}
Halve even \(n\), add \(28\) if \(n\equiv3\pmod{28}\), and add \(7\) to all
other odd \(n\). Which starts diverge? The class \(3\bmod 28\) is a trap
because each subsequent step adds \(28\). Modulo \(7\), both additions
preserve the residue and halving multiplies by \(4\). Thus the three
sets in the first column below are invariant.

Temporarily replace the trap branch by \(n\mapsto n+7\). The resulting map
decreases above \(7\) in at most two steps, since its odd block ends at
\((n+7)/2\). Every orbit reaches \([1,7]\), and direct evaluation of these
seven starts yields the cycles below. A trajectory agrees with this auxiliary
map until the first trap visit. Only the cycle through \(3\) meets the trap,
so exactly the starts in \(\{3,5,6\}\bmod 7\) diverge.
\begin{center}
\small
\begin{tabular}{@{}lll@{}}
\toprule
Residues modulo \(7\)&Auxiliary cycle&Fate with the trap\\
\midrule
\(0\)&\((7,14)\)&Reaches \((7,14)\)\\
\(1,2,4\)&\((1,8,4,2)\)&Reaches \((1,8,4,2)\)\\
\(3,5,6\)&\((3,10,5,12,6)\)&Diverges\\
\bottomrule
\end{tabular}
\end{center}

\newpage
\tier{3}{A binary finite-state classification}
\record{up\_tok3}
Halve even \(n\), add \(3\) if \(n\equiv1\pmod4\), and add \(4\) if
\(n\equiv3\pmod4\). The last class is an increasing trap. Strip the
initial halvings from the orbit. If its odd part is \(3\bmod4\), it diverges.
Otherwise the odd part is \(4B+1\), which reaches \(B+1\) in three steps.

Write \(Q(B)\) for the assertion that \(B+1\) reaches the cycle \((1,4,2)\).
Then \(Q(0)\) holds, \(Q(2Y+1)=Q(Y)\) by halving, and \(Q(4Z) = Q(Z)\)
for \(Z\ge1\) by the same three steps. The input \(4Z+2\) gives a trapped
start \(4Z+3\). The recursions decrease the argument and read the bits of
\(B\) from the low end, consuming a \(1\) or a pair \(00\). An automaton
starts in a state that accepts the end of input, stays there on a \(1\), and
moves on a \(0\) to a state that requires one more \(0\) to return. An
unpaired \(0\) causes rejection. Equivalently the canonical binary expansions
of the original starts that reach the cycle match \texttt{(1(1|00)*01|1)0*}.

The periodic basin contains every power of two. It has density zero, since
for lengths up to \(k\) the expression counts \(O(\varphi^k)\) words, where
\(\varphi=(1+\sqrt5)/2\) comes from tilings by single bits and pairs of bits.
For \(2^{k-1}\le N<2^k\) this bound gives a proportion tending to zero.
An infinite eventually periodic set has positive density. The tier here
reflects the finite-state form of the classification, which follows from
an elementary recursion.

\tier{4}{Steiner's one-minimum cycle theorem}
\record{steiner\_1cycle}
For the shortcut Collatz map, use \(n/2\) on even \(n\) and \((3n+1)/2\)
on odd \(n\). A cycle with one local minimum has one odd run of length
\(k\ge1\) and one even run of length \(\ell\ge1\). If its minimum is
\(2^k b-1\), with \(b\) odd, the odd steps end at \(3^k b-1\) because each
step multiplies \(n+1\) by \(3/2\). Returning to the minimum gives
\[
b\bigl(2^{k+\ell}-3^k\bigr)=2^\ell-1.
\]
Hence \(d=2^{k+\ell}-3^k\) is positive and divides \(2^\ell-1\).
In particular \(2^\ell(2^k-1)<3^k\), so \(\ell\le k\) and
\[
0<\Lambda=(k+\ell)\log 2-k\log 3
=\log(1+d/3^k)<\frac{1}{2^k-1}\le 2^{1-k}.
\]
Rhin's bound, in the form applied by Simons~[2], gives
\(\Lambda\ge(k+\ell)^{-13.3}\ge(2k)^{-13.3}\). This forces \(k\le103\):
at \(k=104\) one has \(208^{133}<2^{1030}\), and
\((k-1)\log 2-13.3\log(2k)\) increases thereafter.
For each \(k\), take the least \(\ell\) with \(d>0\). Larger \(\ell\) violate
\(2^\ell(2^k-1)<3^k\). Testing divisibility for those 103 pairs leaves
only \((k,\ell)=(1,1)\), giving \((1,2)\). Steiner originally proved this
cycle theorem using linear forms in logarithms~[1].

\vspace{4pt}
{\small
\textbf{References.} [1] R. P. Steiner, \emph{A theorem on the Syracuse
problem}, Proc. 7th Manitoba Conf. on Numerical Mathematics and Computing
(1977), 553--559. [2] J. L. Simons, \emph{Cycles and divergent trajectories
for a class of permutation sequences}, Lemma 10,
\href{https://arxiv.org/abs/2205.10582}{arXiv:2205.10582} (2022), applying
G. Rhin's 1987 proposition, p.~160.\par}

\vspace{4pt}
{\small
\textbf{Open catalog.} \texttt{census/}\allowbreak\texttt{open\_collatz\_systems.txt}
covers 66 entries. In CLS-057 (\texttt{qrA3\_q17}), every orbit
reaches a known cycle, and the two basins among nonmultiples of \(17\)
have no eventually periodic residue classification. The open question is
whether a finite automaton can distinguish those basins in some integer base.\par}
\end{document}
